\documentclass[11pt,a4paper]{article}
\usepackage[T1]{fontenc}
\usepackage[utf8]{inputenc}
\usepackage{lmodern}
\usepackage[a4paper,margin=1in]{geometry}
\usepackage{microtype}
\usepackage{enumitem}
\usepackage{amsmath,amssymb}
\usepackage{mathtools}
\usepackage{hyperref}

\usepackage{mathrsfs}

\hypersetup{
	colorlinks=true,
	linkcolor=black,
	urlcolor=blue,
	citecolor=black
}

\setlist[itemize]{leftmargin=1.2em}
\setlist[enumerate]{leftmargin=1.2em}

\newcommand{\Title}[1]{\section*{#1}\vspace{-0.25em}\hrule\vspace{0.8em}}
\newcommand{\Field}[2]{\textbf{#1:}~#2\par}
\newcommand{\Affil}[1]{\textit{#1}\par}
\newcommand{\Abstract}{\vspace{0.6em}\textbf{Abstract.}~}

\begin{document}
	
	\begin{center}
		{\LARGE Online Algebra, Logic, and Topology Seminar}\\[0.25em]
		{\large Abstract}\\[0.5em]
		March 17th, 2026
	\end{center}
	
	\vspace{1.25em}\hrule\vspace{1.25em}
	
	\Title{Dual presentations of locales}
	
	\Field{Speaker}{Graham Manuell}
	\Affil{\hspace{1.15cm}Stellenbosch University, SA}
	
	\noindent \Field{Date \& Time}{March 17th, 2026 (Tuesday) — 3{:}00\,PM (Coimbra time)}
	\noindent \Field{Place}{Google Meet}
	
	\Abstract
Just as usual presentations of frames allow us to define frames using sets of generators and relations, we will discuss how to define frames using compact Hausdorff locales of generators and relations. I will explain how to present a compact Hausdorff locale in terms of itself and describe the locale of partial surjections from Cantor space to a given compact Hausdorff locale. We show the latter is nontrivial, which gives a generalisation of the Alexandroff-Hausdorff theorem without any cardinality restrictions.

\end{document}
